\section{Generación, conservación y realización}
\label{nuclear:introduccion}
El objeto de la Holografía Modular Triádica es un estado generado por posiciones y operaciones que conserva la información de su evolución. La evaluación numérica y la incidencia de frontera son realizaciones correlacionadas de ese antecedente. La contracción de cilindros compatibles determina valores; la prolongación de sus incidencias mantiene relaciones que distinguen los estados generadores. Esta prioridad de la construcción gobierna también la Mecánica Dimensional: acción, duración, extensión y masa aparecen mediante lectores cuyo dominio conserva la orientación, el acarreo y la historia de las operaciones.

El presente artículo estudia una cuestión estructural precisa: qué se conserva cuando una lectura reduce el estado, dónde reside el complemento y cómo reconstruir a partir de él tanto el valor como la dinámica. Una respuesta completa exige tres niveles. Primero se construyen las emisiones y su refinamiento. Después se demuestra una identidad de transporte sobre el estado realizado. Finalmente se determina qué observa un lector concreto y qué consecuencias físicas produce su composición con los lectores de acción y escala. El desarrollo sigue ese orden en cada aplicación.

\subsection{Tesis y enunciados rectores}
La conservación evolutiva de la unidad adquiere una expresión operatoria. Para dos transportes \(B,C\) entre espacios provistos de formas \(b,b'\), con
\[
B^{*}b'B=C^{*}b'C=b,
\]
la representación colectiva de las nueve posiciones produce
\begin{equation}
\mathcal U(B,C)=\frac19
\begin{pmatrix}8B+C&\sqrt8(C-B)\\ \sqrt8(C-B)&B+8C\end{pmatrix},
\qquad
\mathcal U(B,C)^*(b'\oplus b')\mathcal U(B,C)=b\oplus b.
\label{nuclear:eq:rectora}
\end{equation}
La primera columna contiene la lectura y la memoria generadas por un estado sin memoria entrante. La segunda columna describe la contribución de memoria que vuelve a intervenir. La igualdad conserva una forma positiva cuando \(b\) es un producto interior, y conserva área orientada cuando \(b\) es una forma simpléctica. Sus dos interpretaciones utilizan estructuras explícitamente diferentes.

Esta ley permite demostrar cuatro consecuencias conectadas. La primera es la reconstrucción estable del registro completo, incluido su terminal. La segunda recupera de la memoria dodecafásica la dirección que determina el proyector dimensional. La tercera prolonga la conservación a áreas, volúmenes y grados exteriores superiores. La cuarta identifica, en una realización cuántica declarada, la parte de la memoria que altera la pureza de una lectura parcial. Cada consecuencia se obtiene por una operación matemática especificada sobre la anterior.

\begin{figure}[htbp]
\centering
\begin{tikzpicture}[every node/.style={font=\small,align=center},b/.style={draw,rounded corners,inner sep=6pt,text width=39mm},>=Stealth]
\node[b] (g) {Emisión y estado\\APP--TRIT--TPK};
\node[b,right=12mm of g] (u) {Transporte completo\\y forma conservada};
\node[b,right=12mm of u] (l) {Lectura y memoria\\con reconstrucción};
\node[b,below=12mm of g] (i) {Incidencia y\\refinamiento};
\node[b,below=12mm of u] (a) {Acción y\\realización dimensional};
\node[b,below=12mm of l] (s) {Covarianza, espectro\\y entropía reducida};
\draw[->] (g)--(u);\draw[->] (u)--(l);\draw[->] (g)--(i);
\draw[->] (u)--(a);\draw[->] (l)--(s);\draw[->] (i)--(a);\draw[->] (a)--(s);
\end{tikzpicture}
\caption{Dependencia de las realizaciones. Las flechas representan operaciones demostradas en el cuerpo del artículo; la lectura escalar conserva un antecedente más rico que su propio valor.}
\end{figure}

\subsection{Memoria completa y observación parcial}
Para una sucesión de compresiones \(T_n\) con complementos \(D_n\), se escribe
\(P_0=I\), \(P_{n+1}=T_nP_n\). La igualdad local
\(T_n^*T_n+D_n^*D_n=I\) implica
\[
I=A_\infty+\sum_{n\ge0}P_n^*D_n^*D_nP_n,
\qquad
A_\infty=\underset{N\to\infty}{\operatorname{s-lim}}P_N^*P_N.
\]
La reconstrucción utiliza \(A_\infty^{1/2}v\) y todos los \(D_nP_nv\). La convergencia fuerte del Gram terminal basta para la conservación; se demuestra por monotonía positiva, con independencia de la convergencia del vector terminal. El registro contiene amplitudes y correlaciones, mientras sus intensidades proporcionan el balance de norma. Esta distinción determina la información recuperable.

En un registro finito concreto, la incidencia de doce posiciones en las hexadas de Witt posee Gram
\(36I+30\mathbf1\mathbf1^*\). Su normalización conserva exactamente la parte centrada del registro. Al incorporar la media se recuperan sus doce coordenadas. El mapa dimensional posterior utiliza esa información transportada y conserva la dependencia respecto del registro ordenado \(K\).

\subsection{Una consecuencia espectral calculable}
Sea \(J^2=-I\), \(J^{\mathsf T}=-J\), y sea \(H\) simpléctico en una realización canónica de dimensión \(2d\). Se define
\[
D=\frac{\sqrt8}{9}(H-I),\qquad
D_{\mathrm{anti}}=\frac{D+JDJ}{2},\qquad
W=\frac{8I+HH^{\mathsf T}}9.
\]
El artículo demuestra
\begin{equation}
\boxed{-(JW)^2=I+[D,J][D,J]^{\mathsf T}
=I+4D_{\mathrm{anti}}D_{\mathrm{anti}}^{\mathsf T}.}
\label{nuclear:eq:espectro}
\end{equation}
Para dos entradas gaussianas puras idénticas y el transporte \(\mathcal U(I,H)\), la matriz \(W\) es la covarianza reducida normalizada. Por tanto,
\[
\operatorname{Tr}\rho_{\mathrm{vis}}^2
=\det\bigl(I+4D_{\mathrm{anti}}D_{\mathrm{anti}}^{\mathsf T}\bigr)^{-1/4}.
\]
La fórmula relaciona un operador de memoria con el funcional de pureza del estado reducido. La pureza global permanece igual a uno. En particular, la mezcla depende de la componente antilineal respecto de la estructura compleja transportada, y no de una identificación indiscriminada entre memoria y entropía.

La elección de la representación gaussiana, del estado de entrada y del subsistema leído forma parte del enunciado. La genealogía fija los transportes que se componen; el protocolo fija la observación cuya distribución se calcula. Esta especificación permite una comparación reproducible entre realizaciones y separa el teorema de cualquier atribución experimental ulterior.

\subsection{Orden de exposición}
Los capítulos iniciales desarrollan íntegramente las primitivas comunes, la prolongación, las regiones numéricas, el registro y la incidencia excepcional. Siguen los lectores de acción y las representaciones dimensionales. La segunda parte reúne las pruebas de conservación, recuperación, holonomía y extensión reversible. La parte final compone esos resultados con acciones físicas y obtiene el espectro de la memoria reducida. La bibliografía sitúa las herramientas de representación después de la construcción que realizan.

La simetría finita se expresa mediante acciones y entrelazadores. La simetría variacional continua se expresa mediante una acción, una variación y su carga conservada. El teorema de Noether se aplica en este segundo dominio; el transporte de incidencia excepcional conserva su propio contenido algebraico. La relación entre ambos se establece por los observables y las formas transportadas, manteniendo explícitas sus diferencias de tipo.

% BEGIN R27_FRAMING_INTRO_ES
La conservación se examina también antes de introducir representaciones
lineales. La codificación decimal del emisor posee una congruencia de
control y un alfabeto completo de \(42\) bloques, demostrado mediante
sus firmas y semillas. Las incidencias entre mitades de emisiones unen
las regiones conservando la procedencia de cada arista. El contraste
de dos registros con histograma común y correlación diferente establece
una distinción concreta entre distribución y orden. Las leyes de archivo
reciben ese orden junto con las amplitudes y los datos terminales
que figuran en sus enunciados.

La normalización del logaritmo multiplicativo se distingue de su
inversión mediante el criterio del menor generador positivo. Sobre la
clausura entera, los controles de sentido de acarreo, extremo terminal,
rotación y perturbación local mantienen fijos los restantes datos.
La selección de frontera aplica sucesivamente saturación, neutralidad
dual y orientación. Estas operaciones preceden a las lecturas de
recuperación y conservan su dominio, de modo que cada inversa actúa
sobre los datos efectivamente producidos por su construcción.
% END R27_FRAMING_INTRO_ES

% BEGIN R28_FRAMING_INTRO_ES
El enlace de los registros se explicita mediante el multigrafo del
emisor y su cociente de fase: cinco cierres mínimos reúnen las cuatro
incidencias requeridas para cada región. La derivación de la ocupación
de frontera y la paridad de los lectores bajo inversión orientada
completan la descripción de los datos que recibe el transporte bilateral
y que mantienen sus lecturas de acoplamiento.
% END R28_FRAMING_INTRO_ES
